\documentclass[letterpaper]{article}


\newcommand{\contribution}[1]{\iftoggle{contrib}{\emph{Contribution: #1}}}

\usepackage[hidelinks]{hyperref}
\usepackage{url}
\usepackage{breakurl}
\usepackage{geometry}


\usepackage{etoolbox}




% Comment the following lines to use the default Computer Modern font
% instead of the Palatino font provided by the mathpazo package.
% Remove the 'osf' bit if you don't like the old style figures.
\usepackage[T1]{fontenc}
\usepackage[sc,osf]{mathpazo}

% Set your name here
\def\name{Dr. Casey William Dunn}

% Replace this with a link to your CV if you like, or set it empty
% (as in \def\footerlink{}) to remove the link in the footer:
\def\footerlink{}

% The following metadata will show up in the PDF properties
%\hypersetup{
%  colorlinks = true,
%  urlcolor = black,
%  pdfauthor = {\name},
%  pdfkeywords = {evolution, invertebrate zoology, computational biology},
%  pdftitle = {\name: Curriculum Vitae},
%  pdfsubject = {Curriculum Vitae},
%  pdfpagemode = UseNone
%}

\geometry{
  body={6.5in, 8.5in},
  left=1.0in,
  top=1.25in
}

% Customize page headers
\pagestyle{myheadings}
\markright{\name}
\thispagestyle{empty}

% Custom section fonts
\usepackage{sectsty}
\sectionfont{\rmfamily\mdseries\Large}
\subsectionfont{\rmfamily\mdseries\itshape\large}

% Other possible font commands include:
% \ttfamily for teletype,
% \sffamily for sans serif,
% \bfseries for bold,
% \scshape for small caps,
% \normalsize, \large, \Large, \LARGE sizes.

% Don't indent paragraphs.
\setlength\parindent{0em}

% Make lists without bullets
\renewenvironment{itemize}{
  \begin{list}{}{
    \setlength{\leftmargin}{1.5em}
  }
}{
  \end{list}
}


\newcommand{\entrylabel}[1]{\mbox{{\rm #1}}\hfill}
\newenvironment{entry}[1]
{\begin{list}{}%
{\renewcommand{\makelabel} {\entrylabel}%
\setlength{\labelwidth}{#1}%
\setlength{\leftmargin}{#1}%
\addtolength{\leftmargin}{2mm}
\setlength{\parsep}{0pt}
}%
}%
{\end{list}}


\begin{document}

% Toggles for conditional inclusion:
% http://tex.stackexchange.com/questions/5894/latex-conditional-expression/5896#5896

% Section toggles



% Show home address
\newtoggle{home}
\togglefalse{home}
% \toggletrue{home}

% Individual contributions to papers
\newtoggle{contrib}
\togglefalse{contrib}
% \toggletrue{contrib}

\newtoggle{inreview}
\togglefalse{inreview}
% \toggletrue{inreview}

\newtoggle{inprep}
\togglefalse{inprep}
% \toggletrue{inprep}

\newtoggle{grants}
\togglefalse{grants}
% \toggletrue{grants}

\newtoggle{declinedgrants}
\togglefalse{declinedgrants}
% \toggletrue{declinedgrants}

\newtoggle{service}
\togglefalse{service}
% \toggletrue{service}

\newtoggle{long}
\toggletrue{long}
% \toggletrue{service}

% Master toggles, default is minimal public (web) view

% Outside review
% \toggletrue{grants} \toggletrue{service}

% Yale internal
%\toggletrue{grants} \toggletrue{service} \toggletrue{inreview}  %\toggletrue{inprep}

% Promotion/ TPAC
% \toggletrue{grants} \toggletrue{service} \toggletrue{contrib} \toggletrue{inreview} \toggletrue{inprep}

% Short
% \togglefalse{long}  \toggletrue{grants}



% To get bib entry and hyper ref to work together...
%\begingroup
%\makeatletter
%\let\@bibitem\saved@bibitem
%%\nobibliography{database}
%\nobibliography{Casey_Dunn_CV.bib}
%\bibliographystyle{apalike} %be careful here, there is only a few styles that will run
%\endgroup


%\nobibliography{Casey_Dunn_CV.bib}
%\bibliographystyle{apalike} %be careful here, there is only a few styles that will run



% Place name at left
\noindent {\huge \name -- Curriculum Vitae}

% Alternatively, print name centered and bold:
%\centerline{\huge \bf \name}

%\vspace{0.25in}

\section{Position and Contact Information}

\begin{minipage}[b]{0.45\linewidth}
  Professor \\ 

  Ecology and Evolutionary Biology \\
  Yale University \\
  165 Prospect St. \\
  New Haven, CT 06511 USA\\
    casey.dunn@yale.edu \\

\end{minipage}
\begin{minipage}[b]{0.45\linewidth}

    \url{http://dunnlab.org/} \\
    \url{https://github.com/caseywdunn} \\
    \url{https://scholar.google.com/citations?user=mSMSMM4AAAAJ} \\
    \url{https://orcid.org/0000-0003-0628-5150} \\
\end{minipage}


\iftoggle{home}{

\section{Home address}

New Haven, CT 06511%\\
%USA

}% end \iftoggle{home}

\section{Education}

\begin{entry}{60pt}
 \item [\rm 2005] Ph.D. Yale University, Department of Ecology and Evolutionary Biology (With Distinction)
 \item [\rm 2003] M.S. Yale University, Department of Ecology and Evolutionary Biology
 \item [\rm 1999] B.S. Stanford University, Biological Sciences (Honors, Phi Beta Kappa)
\end{entry}


\section{Professional Appointments}

%\begin{tabular}{ll}
\begin{entry}{60pt}
\item [\rm 2021--present] Curator, Informatics, Peabody Museum, Yale University
\item [\rm 2017--present] Professor, Ecology and Evolutionary Biology, Yale University
\item [\rm 2019--present] Professor (Secondary), Genetics, Yale University
\item [\rm 2017--present] Curator, Division of Invertebrate Zoology, Peabody Museum, Yale University
\item [\rm 2017--2018] Adjunct Professor, Ecology and Evolutionary Biology, Brown University
\item [\rm 2014--2017]  Associate Professor, Ecology and Evolutionary Biology, Brown University
\iftoggle{long}{
\item [\rm 2010--2014]  Manning Assistant Professor, Ecology and Evolutionary Biology, Brown University
\item [\rm 2008--2010] Assistant Professor, Ecology and Evolutionary Biology, Brown University
}{
\item [\rm 2008--2014] Assistant Professor, Ecology and Evolutionary Biology, Brown University
} %\iftoggle{long}
\item [\rm 2007--2008] Assistant Professor (Research), Ecology and Evolutionary Biology, Brown University
\item [\rm 2005--2007] Postdoctoral Fellow, Kewalo Marine Lab, University of Hawaii
\item [\rm 1999--2000] Built and maintained eddy covariance systems in the lab of Stuart Chapin III at the University of Alaska Fairbanks, under the supervision of James Randerson of Caltech 
\item [\rm 1998--1999] Maintained and built cryogenic research equipment as student research assistant with Doug Osheroff, Physics Department, Stanford University
\item [\rm 1995--1998] Co-op Technical at KLA-Tencor. Built equipment and software to monitor instruments used in semiconductor fabrication
\item [\rm 1994--1995] Tour guide at La Selva Jungle Lodge in the Amazon Basin of eastern Ecuador
\item [\rm 1993--1994] Student researcher in John Postlethwait's lab, University of Oregon
\end{entry}

\iftoggle{long}{
\section{Publications}
}{
\section{Selected Publications}
} %\iftoggle{long}


\subsection*{a. Books/ monographs}

\begin{itemize}

\item SHD Haddock and CW Dunn. (2010) Practical Computing for Biologists. 538 pages. Sinauer Associates. Sunderland, MA. \url{http://practicalcomputing.org}. \contribution{My co-author and I participated equally in all phases of writing this book, from conceptualization through writing and editing. We collaborated on all chapters. We determined the author order with a random number generator, as described on page 220 of the book.}

\end{itemize}

\iftoggle{long}{

\subsection*{b. Chapters in books}

\begin{itemize}

\item Giribet, G, CW Dunn, GD Edgecombe, A Hejnol, MQ Martindale, and GW Rouse (2009) Assembling the spiralian tree of life (p 52-64) In: Animal Evolution. MJ Telford and DTJ Littlewood (eds). Oxford University Press. \contribution{I contributed text about phylogenomic analyses, and helped integrate and revise the manuscript.}

\item Mills, CE, AC Marques, AE Migotto, DR Calder, C Hand, JT Rees, SHD Haddock, CW Dunn, and PR Pugh. (2007) Hydrozoa: Polyps, Hydromedusae, and Siphonophora. In: The Light \& Smith manual: intertidal invertebrates from central California to Oregon, 4th edition. JT Carlton (ed). University of California Press. \contribution{I helped write the siphonophore key.}

\end{itemize}

} % end \iftoggle{long}

\subsection*{c. Refereed journal articles}

\begin{itemize}

\item Oguchi K, Maeno A, Yoshida K, Kohtsuka H, Dunn CW (2025) Zooid arrangement and colony growth in Porpita porpita. Frontiers in Zoology 22:11. \url{https://doi.org/10.1186/s12983-025-00565-3}

\item Church SH, Abedon RB, Ahuja N, Anthony CJ, Destanović D, Ramirez DA, Rojas LM, Albinsson ME, Álvarez Trasobares I, Bergemann RE, Bogdanovic O, Burdick DR, Cunha TJ, Damian-Serrano A, D’Elía G, Dion KB, Doyle TK, Gonçalves JM, Gonzalez Rajal A, Haddock SHD, Helm RR, Le Gouvello D, Lewis ZR, Magalhães BIMM, Mańko MK, Mayorga-Adame CG, de Mendoza A, Moura CJ, Munro C, Nel R, Oguchi K, Perelman JN, Prieto L, Pitt KA, Roughan M, Schaeffer A, Schmidt AL, Sellanes J, Wilson NG, Yamamoto G, Lazo-Wasem EA, Simon C, Decker MB, Coughlan JM, Dunn CW (2025) Population genomics of a sailing siphonophore reveals genetic structure in the open ocean. Current Biology. \url{https://doi.org/10.1016/j.cub.2025.05.066}

\item Kantor R, Steingrimsson J, Fulton J, Novitsky V, Howison M, Gillani F, Bhattarai L, MacAskill M, Hague J, Guang A, Khanna A, Dunn C, Hogan J, Bertrand T, Bandy U (2024) Prospective evaluation of routine statewide integration of molecular epidemiology and contact tracing to disrupt HIV transmission. Open Forum Infectious Diseases. \url{https://doi.org/10.1093/ofid/ofae599}

\item Oguchi K, Yamamoto G, Kohtsuka H, Dunn CW (2024) Physalia gonodendra are not yet sexually mature when released. Scientific Reports 14:23011. \url{https://doi.org/10.1038/s41598-024-73611-5}

\item Fulton J, Novitsky V, Gillani F, Guang A, Steingrimsson J, Khanna A, Hague J, Dunn C, Hogan J, Howe K, MacAskill M, Bhattarai L, Bertrand T, Bandy U, Kantor R (2024) Integrating HIV Cluster Analysis in Everyday Public Health Practice: Lessons Learned From a Public Health–Academic Partnership. JAIDS Journal of Acquired Immune Deficiency Syndromes 97(1):48-54. \url{https://doi.org/10.1097/QAI.0000000000003469}

\item Church SH, Mah JL, Dunn CW (2024) Integrating phylogenies into single-cell RNA sequencing analysis allows comparisons across species, genes, and cells. PLOS Biology 22(5): e3002633. \url{https://doi.org/10.1371/journal.pbio.3002633}

\item N Ahuja, X Cao, DT Schultz, N Picciani, A Lord, S Shao, K Jia, DR Burdick, SH D Haddock, Y Li, CW Dunn (2024) Giants among Cnidaria: Large Nuclear Genomes and Rearranged Mitochondrial Genomes in Siphonophores. Genome Biology and Evolution, 16(3). \url{https://doi.org/10.1093/gbe/evae048}

\item ED Hetherington, HG Close, SH Haddock, AD Damian-Serrano, CW Dunn, NJ Wallsgrove, SC Doherty, CA Choy (2024) Vertical trophic structure and niche partitioning of gelatinous predators in a pelagic food web: Insights from stable isotopes of siphonophores. Limnology and Oceanography. \url{https://doi.org/10.1002/lno.12536}

\item JL Mah, CW Dunn (2024) Cell type evolution reconstruction across species through cell phylogenies of single-cell RNA sequencing data. Nature Ecology and Evolution. \url{https://doi.org/10.1038/s41559-023-02281-9}

\item A Khanna, V Novitsky, A Guang, M Howison, FS Gillani, J Steingrimsson, C Dunn, J Fulton, T Bertrand, K Howe, L Bhattarai, G Ronquillo, M MacAskill, U Bandy, R Kantor (2023) Integrating HIV Partner Services and Molecular Epidemiology Data to Enhance HIV Transmission Disruption in Rhode Island: Findings from a Public Health-Academic Partnership. Open Forum Infectious Diseases, 10(2).  \url{https://doi.org/10.1093/ofid/ofad500.170}

\item SH Church, JL May, G Wagner, CW Dunn (2023) Normalizing need not be the norm: count‑based math for analyzing single‑cell data. Theory in Biosciences. \url{https://doi.org/10.1007/s12064-023-00408-x}

\item M Howison, FS Gillani, V Novitsky, JA Steingrimsson, J Fulton, T Bertrand, K Howe, A Civitarese, L Bhattarai, M MacAskill, G Ronquillo, J Hague, CW Dunn, U Bandy, JW Hogan, R Kantor (2023) An Automated Bioinformatics Pipeline Informing Near-Real-Time Public Health Responses to New HIV Diagnoses in a Statewide HIV Epidemic. Viruses, 15(3), 737. \url{https://www.mdpi.com/1999-4915/15/3/737}

\item CW Dunn (2023) Neurons that connect without synapses. Science, Vol 380, Issue 6642, pp. 241-242. \url{https://doi.org/10.1126/science.adh0542}

\item SH Church, C Munro, CW Dunn, CG Extavour (2023) The evolution of ovary-biased gene expression in Hawaiian Drosophila. PLoS Genet 19(1): e1010607. \url{https://doi.org/10.1371/journal.pgen.1010607}

\item V Novitsky, J Steingrimsson, M Howison, CW Dunn, FS Gillani, J Fulton, T Bertrand, K Howe, L Bhattarai, G Ronquillo, M MacAskill, U Bandy, J Hogan, R Kantor (2023) Not all clusters are equal: dynamics of molecular HIV-1 clusters in a statewide Rhode Island epidemic. AIDS. \url{https://doi.org/10.1097/QAD.0000000000003426}

\item A Monroe-Wise, JA Steingrimsson, J Fulton, M Howison, V Novitsky, FS Gillani, T Bertrand, A Civitarese, K Howe, G Ronquillo, B Lafazia, Z Parillo, T Marak, PA Chan, L Bhattarai, CW Dunn, U Bandy, NA Scott, JW Hogan, R Kantor (2022) Beyond HIV outbreaks: protocol, rationale and implementation of a prospective study quantifying the benefit of incorporating viral sequence clustering analysis into routine public health interventions. BMJ Open, 12(4). \url{https://bmjopen.bmj.com/content/12/4/e060184}

\item A Guang, M Howison, L Ledingham, M D’Antuono, PA Chan, C Lawrence, CW Dunn, R Kantor (2022) Incorporating Within-Host Diversity in Phylogenetic Analyses for Detecting Clusters of New HIV Diagnoses. Frontiers in Microbiology. \url{https://www.frontiersin.org/articles/10.3389/fmicb.2021.803190/full}

\item JA Steingrimsson, J Fulton, MH Howison, V Novitsky, FS Gillani, T Bertrand, A Civitarese, K Howe, G Ronquillo, B Lafazia, Z Parillo, T Marak, PA Chan, L Bhattarai, CW Dunn, U Bandy, NA Scott, JW Hogan, R Kantor (2022) Beyond HIV outbreaks:  protocol, rationale, and implementation of a prospective study quantifying the benefit of incorporating viral sequence clustering analysis into routine public health interventions. BMJ Open. \url{http://dx.doi.org/10.1136/bmjopen-2021-060184}

\item Damian-Serrano A, Hetherington ED, Choy A, Haddock SHD, Lepides A, Dunn CW (2022) Characterizing the secret diets of siphonophores (Cnidaria: Hydrozoa) using DNA metabarcoding. PLoS One. \url{https://doi.org/10.1371/journal.pone.0267761}. GitHub repository \url{https://github.com/dunnlab/siphweb_metabarcoding}.

\item Hetherington ED, Damian-Serrano A, Haddock SHD, Dunn CW, Choy A (2022) Integrating siphonophores into marine food-web ecology. Limnology and Oceanography Letters. \url{https://doi.org/10.1002/lol2.10235}

\item Munro C, Zapata F, Howison M, Siebert S, Dunn CW (2022) Evolution of gene expression across species and specialized zooids in Siphonophora. MBE. \url{https://doi.org/10.1093/molbev/msac027}

\item Li, Y, Shen XX, Evans B, Dunn CW, Rokas A. (2021) Rooting the animal tree of life. MBE. \url{https://doi.org/10.1093/molbev/msab170}

\item Novitsky, V, Steingrimsson J, Howison M, Dunn CW, Gillani FS, Manne A, Li Y, Spence M, Parillo Z, Fulton J, Marak T, Chan P, Bertrand T, Bandy U, Alexander-Scott N, Hogan J, Kantor R. (2021) Longitudinal typing of molecular HIV clusters in a statewide epidemic. AIDS. \url{https://doi.org/10.1097/QAD.0000000000002953}

\item Damian-Serrano, A, SHD Haddock, CW Dunn (2021) The evolutionary history of siphonophore tentilla: Novelties, convergence, and integration. Integrative Organismal Biology. \url{https://doi.org/10.1093/iob/obab019}

\item Damian-Serrano, A, SHD Haddock, CW Dunn (2021) The evolution of siphonophore tentilla for specialized prey capture in the open ocean. PNAS 118 (8) e2005063118. \url{https://doi.org/10.1073/pnas.2005063118} (Cover article)

\item Li Y, Steenwyk JL, Chang Y, Wang Y, James TY, Stajich JE, Spatafora JW, Groenewald M, Dunn CW, Hittinger CT, Shen XX, Rokas A (2021) A genome-scale phylogeny of the kingdom Fungi. Current Biology 31:8 1653-1665. \url{https://doi.org/10.1016/j.cub.2021.01.074}

\item Guang, A, M Howison, F Zapata, C Lawrence, CW Dunn (2021) Revising transcriptome assemblies with phylogenetic information. PLoS ONE 16(1): e0244202.  \url{https://doi.org/10.1371/journal.pone.0244202}. Git code repository: \url{https://github.com/caseywdunn/ms_treeinform}


\item Kantor, R, JP Fulton, J Steingrimsson, V Novitsky, M Howison, F Gillani, Y Li, A Manne, Z Parillo , M Spence, T Marak, P Chan, CW Dunn, T Bertrand, U Bandy, N Alexander-Scott, JW Hogan (2020) Challenges in evaluating the use of viral sequence data to identify HIV transmission networks for public health. Statistical Communications in Infectious Diseases 12:s1. \url{https://doi.org/10.1515/scid-2019-0019}


\item Novitsky, V, JA Steingrimsson, M Howison, FS Gillani, Y Li, A Manne, J Fulton, M Spence, Z Parillo, T Marak, PA Chan, T Bertrand, U Bandy, N Alexander-Scott, CW Dunn, J Hogan, R Kantor. (2020) Empirical comparison of analytical approaches for identifying molecular HIV-1 clusters. Scientific Reports 10, Article number: 18547. \url{https://doi.org/10.1038/s41598-020-75560-1}

\item AG Collins, M Daly, and CW Dunn. 2020. Cnidaria. Pp. 457–460 in Phylonyms: A Companion to the PhyloCode (K. de Queiroz, P. D. Cantino, and J. A. Gauthier, eds.). CRC Press, Boca Raton, FL. \url{https://doi.org/10.1201/9780429446276}

\item AG Collins and CW Dunn. 2020. Medusozoa. Pp. 473-475 in Phylonyms: A Companion to the PhyloCode (K. de Queiroz, P. D. Cantino, and J. A. Gauthier, eds.). CRC Press, Boca Raton, FL. \url{https://doi.org/10.1201/9780429446276}

\item Dunn, CW and AG Collins. 2020. Hydrozoa. Pp. 473-475 in Phylonyms: A Companion to the PhyloCode (K. de Queiroz, P. D. Cantino, and J. A. Gauthier, eds.). CRC Press, Boca Raton, FL. \url{https://doi.org/10.1201/9780429446276}

\item Dunn, CW and AG Collins. 2020. Trachylina. Pp. 485-487 in Phylonyms: A Companion to the PhyloCode (K. de Queiroz, P. D. Cantino, and J. A. Gauthier, eds.). CRC Press, Boca Raton, FL. \url{https://doi.org/10.1201/9780429446276}

\item Dunn, CW and AG Collins. 2020. Hydroidolina. Pp. 489-490 in Phylonyms: A Companion to the PhyloCode (K. de Queiroz, P. D. Cantino, and J. A. Gauthier, eds.). CRC Press, Boca Raton, FL. \url{https://doi.org/10.1201/9780429446276}

\item Dunn, CW 2020. Siphonophora. Pp. 491-492 in Phylonyms: A Companion to the PhyloCode (K. de Queiroz, P. D. Cantino, and J. A. Gauthier, eds.). CRC Press, Boca Raton, FL. \url{https://doi.org/10.1201/9780429446276}

\item Smith, SD, MW Pennell, CW Dunn, SV Edwards (2020) Phylogenetics is the New Genetics (for Most of Biodiversity). Trends in Ecology and Evolution. \url{https://doi.org/10.1016/j.tree.2020.01.005}.

\item Catriona Munro, Zer Vue, Richard R. Behringer, Casey W. Dunn (2019) Morphology and development of the Portuguese man of war, Physalia physalis. Scientific Reports 9:15522. \url{http://dx.doi.org/10.1038/s41598-019-51842-1}.

\item Li-Gen Wang, Tommy Tsan-Yuk Lam, Shuangbin Xu, Zehan Dai, Lang Zhou, Tingze Feng, Pingfan Guo, Casey W Dunn, Bradley R Jones, Tyler Bradley, Huachen Zhu, Yi Guan, Yong Jiang, Guangchuang Yu. (2019) treeio: an R package for phylogenetic tree input and output with richly annotated and associated data. Molecular Biology and Evolution. \url{https://doi.org/10.1093/molbev/msz240}. Git code repository: \url{https://guangchuangyu.github.io/software/treeio/}.

\item Munro, C, S Siebert, F Zapata, M Howison, A Damian-Serrano, SH Church, FE Goetz, PR Pugh, SHD Haddock, CW Dunn (2018) Improved phylogenetic resolution within Siphonophora (Cnidaria) with implications for trait evolution. Molecular Phylogenetics and Evolution. \url{http://dx.doi.org/10.1016/j.ympev.2018.06.030}. bioRxiv preprint: \url{https://doi.org/10.1101/251116}. Git code repository: \url{https://github.com/caseywdunn/siphonophore_phylogeny_2017}.

\item Pugh, PR, CW Dunn, SHD Haddock (2018) Description of Tottonophyes enigmatica gen. nov., sp. nov. (Hydrozoa, Siphonophora, Calycophorae), with a reappraisal of the function and homology of nectophoral canals. Zootaxa 4415(3):452-472. \url{https://doi.org/10.11646/zootaxa.4415.3.3}.

\item Haddock, SHD, Christianson LM, Francis WR, Martini S, Dunn CW, Pugh PR, Mills CE, Osborn KJ, Seibel BA, Choy CA, Schnitzler CE, Matsumoto GI, Messi M, Schultz DT, Winnikoff JR, Powers ML, Gasca R, Browne WE, Johnsen S, Schlining KL, von Thun S, Erwin BE, Ryan JF, Thuesen EV (2018) Insights into the Biodiversity, Behavior, and Bioluminescence of Deep-Sea Organisms Using Molecular and Maritime Technology. Oceanography 30:38-47. \url{https://doi.org/10.5670/oceanog.2017.422}

\item Dunn, CW, F Zapata, C Munro, S Siebert, A Hejnol (2018) Pairwise comparisons across species are problematic when analyzing functional genomic data. PNAS. \url{http://dx.doi.org/10.1073/pnas.1707515115}. bioRxiv preprint: \url{https://doi.org/10.1101/107177}. Git code repository: \url{https://github.com/caseywdunn/comparative_expression_2017}.

\item Helm, RR, CW Dunn (2017) Indoles induce metamorphosis in a broad diversity of jellyfish, but not in a crown jelly (Coronatae). PLoS One. \url{https://doi.org/10.1371/journal.pone.0188601}.

\item Dunn, CW, C Munro (2016) Comparative genomics and the diversity of life. Zoologica Scripta 45:5-13. \url{http://dx.doi.org/10.1111/zsc.12211}.

\item Guang, A, F Zapata, M Howison, CE Lawrence, CW Dunn (2016) An Integrated Perspective on Phylogenetic Workflows. Trends in Ecology and Evolution 31:116-126. \url{http://dx.doi.org/10.1016/j.tree.2015.12.007}. \contribution{This paper is a component of Guang's thesis work in my group. It was motivated by a desire to identify statistically principled solutions to engineering challenges and biological questions in the lab. I worked closely with Guang, a mathematician, on identifying our biological goals and exploring possible approaches. I contributed much of the text.}

\item Dunn, CW (2015) Acorn worms in a nutshell. Nature 527:448-449. \url{http://dx.doi.org/10.1038/nature16315}.

\item Zapata, F, FE Goetz, SA Smith, M Howison, S Siebert, S Church, SM Sanders, CL Ames, CS McFadden, SC France, M Daly, AG Collins, SHD Haddock, CW Dunn, P Cartwright (2015) Phylogenomic analyses support traditional relationships within Cnidaria. PLoS One 10(10): e0139068. \url{http://dx.doi.org/10.1371/journal.pone.0139068}. BioRxiv preprint: \url{http://dx.doi.org/10.1101/017632}. Git code repository: \url{https://bitbucket.org/caseywdunn/cnidaria2014}.

\item Dunn, CW, JT Ryan (2015) The evolution of animal genomes. Current Opinion in Genetics and Development 35:25-32. \url{http://dx.doi.org/10.1016/j.gde.2015.08.006}

\item Haddock, SH, CW Dunn (2015) Fluorescent proteins function as a prey attractant: experimental evidence from the hydromedusa Olindias formosus and other marine organisms. Biology Open 4:1094-1104. \url{http://dx.doi.org/10.1242/bio.012138}.

\item Church, SH, JF Ryan, CW Dunn (2015) Automation and Evaluation of the SOWH Test of Phylogenetic Topologies with SOWHAT. Systematic Biology 64(6):1048-1058. \url{http://dx.doi.org/10.1093/sysbio/syv055}. BioRxiv preprint: \url{http://dx.doi.org/10.1101/005264}. Git code repository: \url{https://github.com/josephryan/sowhat}. \contribution{This manuscript began as a student project by Sam Church in my Phylogenetic Biology class. I helped plan and interpret the analyses.}

\item Laumer, CE, N Bekkouche, A Kerbl, F Goetz, RC Neves, MV Sorensen, RM Kristensen, A Hejnol, CW Dunn, G Giribet, K Worsaae (2015) Spiralian Phylogeny Informs the Evolution of Microscopic Lineages. Current Biology 25(15):2000-2006.  \url{http://dx.doi.org/10.1016/j.cub.2015.06.068}.

\item Church, SH, S Siebert, P Bhattacharyya, CW Dunn (2015) The Histology of \emph{Nanomia bijuga} (Hydrozoa: Siphonophora). J. Exp. Zool. (Mol. Dev. Evol.) 324(5):435-449.  \url{http://dx.doi.org/10.1002/jez.b.22629}. BioRxiv preprint: \url{http://dx.doi.org/10.1101/010868}.

\item Siebert, S, FE Goetz, SH Church, P Bhattacharyya, F Zapata, SHD Haddock, and CW Dunn (2015) Stem Cells in a Colonial Animal with Localized Growth Zones. EvoDevo 6:22. \url{http://dx.doi.org/10.1186/s13227-015-0018-2}. BioRxiv preprint: \url{http://dx.doi.org/10.1101/001685}. Git code repository: \url{https://bitbucket.org/caseywdunn/siebert_etal}. \contribution{This manuscript describes the first stem cells to be found in siphonophores, and their relevance to understanding the evolution of stem cells in multi-cellular organisms. All work has been done in my lab, with generous help from collaborators to collect the specimens at sea. Undergraduates played a critical role in this project, with some of the key discoveries coming directly from their observations. I have been involved in all stages of project design and data interpretation.}

\item Salinas-Saavedra, M, TQ Stephenson, CW Dunn, MQ Martindale (2015) Par system components are asymmetrically localized in ectodermal epithelia, but not during early development in the sea anemone \emph{Nematostella vectensis}. EvoDevo 6:20. \url{http://dx.doi.org/10.1186/s13227-015-0014-6}.

\item Helm, RR, S Tiozzo, MKS Lilley, F Lombard, CW Dunn (2015) Comparative muscle development of scyphozoan jellyfish with simple and complex life cycles. EvoDevo 6:11. \url{http://dx.doi.org/10.1186/s13227-015-0005-7}.

\item Dunn, CW, SP Leys, SHD Haddock (2015) The hidden biology of sponges and ctenophores. Trends in Ecology and Evolution 30:282-291. \url{http://dx.doi.org/10.1016/j.tree.2015.03.003}. \contribution{I initiated this manuscript and took the lead on writing it. }

\iftoggle{long}{
\item Reich, A, CW Dunn, K Akasaka, G Wessel (2015) Phylogenomic Analyses of Echinodermata Support the Sister Groups of Asterozoa and Echinozoa. PLoS One 10:e0119627. \url{http://dx.doi.org/10.1371/journal.pone.0119627}.

\item Gonzalez, VL, SCS Andrade, R Bieler, TM Collins, CW Dunn, PM Mikkelsen, JD Taylor, G Giribet (2015) A phylogenetic backbone for Bivalvia: an RNA-seq approach. Proceedings of the Royal Society B: Biological Sciences 282:20142332. \url{http://dx.doi.org/10.1098/rspb.2014.2332}.
}

\item Dunn, CW, G Giribet, GD Edgecombe, A Hejnol (2014) Animal Phylogeny and its Evolutionary Implications. Annual Review of Ecology, Evolution, and Systematics 45:371-395. \url{http://dx.doi.org/10.1146/annurev-ecolsys-120213-091627}. \contribution{I was invited to write this review on animal evolution. Several colleagues have agreed to join me on the manuscript. }

\item Zapata, F, NG Wilson, M Howison, SCS Andrade, KM J\"orger, Michael Schr\"odl, Freya E Goetz, Gonzalo Giribet, Casey W Dunn (2014) Phylogenomic analyses of deep gastropod relationships reject Orthogastropoda. Proceedings of the Royal Society B: Biological Sciences 281:1471-2954. \url{http://dx.doi.org/10.1098/rspb.2014.1739}. BioRxiv preprint: \url{http://dx.doi.org/10.1101/007039}. Git code repository: \url{https://bitbucket.org/caseywdunn/gastropoda}. \contribution{I helped design the study, oversaw most of the data acquisition, helped design the data analysis strategy, and played a large role in writing the manuscript.}

\item Dunn, CW (2014) Reconsidering the phylogenetic utility of miRNA in animals. PNAS 111:12576-12577. \url{http://dx.doi.org/10.1073/pnas.1413545111}.

\item Howison, M, F Zapata, EJ Edwards, and CW Dunn (2014) Bayesian genome assembly and assessment by Markov Chain Monte Carlo sampling.  PLoS One 9:e99497. \url{http://dx.doi.org/10.1371/journal.pone.0099497}. arXiv preprint: \url{http://arxiv.org/abs/1308.1388}. Git code repository: \url{https://bitbucket.org/mhowison/gabi}. Example analysis report: \url{https://web3.ccv.brown.edu/mhowison/gabi-report/}. \contribution{I conceived of the idea of using MCMC to approximate the posterior probabilities of assemblies, provided guidance on implementation (which was written entirely by M Howison), and helped to author the manuscript. }

\iftoggle{long}{
\item Lopez, J, H Bracken-Grissom, A Collins, T Collins, K Crandall, D Distel, C Dunn, et al. (2014) Global Invertebrate Genomics Alliance (GIGA): Developing Community Resources to Study Diverse Invertebrate Genomes. Journal of Heredity 105:1-18. \url{http://dx.doi.org/10.1093/jhered/est084} \contribution{This is a large collaborative effort to coordinate the sequencing of invertebrate genomes. My primary contributions were to help formulate standards for collecting and processing samples, generate sequence data, and perform analyses, as well as to select species for sequencing.}
}

\item Ryan, JF, K Pang, CE Schnitzler, A Nguyen, RT Moreland, DK Simmons, BJ Koch, WR Francis, P Havlak, NISC Comparative Sequencing Program, SA Smith, NH Putnam, SHD Haddock, CW Dunn, TG Wolfsberg, JC Mullikin, MQ Martindale, AD Baxevanis (2013) The genome of the ctenophore \emph{Mnemiopsis leidyi} and its implications for cell type evolution. Science 432:1242592. \url{http://dx.doi.org/10.1126/science.1242592}. \contribution{I helped to design and implement some of the phylogenetic analyses.}

\item CW Dunn, M Howison, and F Zapata (2013) Agalma: an automated phylogenomics workflow. BMC Bioinformatics 14:330. \url{http://dx.doi.org/10.1186/1471-2105-14-330}. arXiv preprint: \url{http://arxiv.org/abs/1307.6432}. Git code repository: \url{https://bitbucket.org/caseywdunn/agalma} (software), \url{https://bitbucket.org/caseywdunn/dunnhowisonzapata2013} (analyses).
\contribution{This manuscript describes the software tool Agalma developed in my lab. Agalma has already been released and is in use by multiple labs. This project has been developed entirely in my lab, and I have been involved in all aspects including methods design, software development, testing, and writing of the manuscript.}

\item Howison, M, F Zapata, CW Dunn (2013) Toward a statistically explicit understanding of \emph{de novo} sequence assembly. Bioinformatics 29:2959-2963. \url{http://dx.doi.org/10.1093/bioinformatics/btt525}. \contribution{This review summarizes current challenges in the representation of uncertainty in genome assemblies, and proposes several new approaches. It is a collaborative effort between myself and a postdoc and application scientist in my lab. All authors took part in the discussions that led to the manuscript, the development of new tools, and the writing of the manuscript.}

\item Siebert, S, PR Pugh, SHD Haddock, CW Dunn (2013) Re-evaluation of characters in Apolemiidae (Siphonophora), with description of two new species from Monterey Bay, California. Zootaxa 3702: 201-232.  \url{http://www.mapress.com/zootaxa/2013/f/zt03702p232.pdf}. Accompanying video of new species available at \url{https://vimeo.com/59928656}. \contribution{This project was led by my postdoc S. Siebert. Most of the work has been done in my lab. I helped design the project, interpret the data, and write the manuscript.}

\item Dunn, CW, X Luo, Z Wu (2013) Phylogenetic analysis of gene expression. Integrative and Comparative Biology 53:847-856. \url{http://dx.doi.org/10.1093/icb/ict068}. arXiv preprint: \url{http://arxiv.org/abs/1302.2978}. Git code repository: \url{https://bitbucket.org/caseywdunn/sicb2013}.
\contribution{I was invited to submit this manuscript to a special issue of the journal resulting from a symposium I participated in. I wrote most of the manuscript and formalized the problem, and collaborated on writing the software and conducting analyses.}

\item Helm, RR, S Siebert, S Tulin, J Smith, CW Dunn (2013) Characterization of differential transcript abundance through time during \emph{Nematostella vectensis} development. BMC Genomics 14:266. \url{http://dx.doi.org/10.1186/1471-2164-14-266}. Git code repository: \url{https://bitbucket.org/caseywdunn/helm_etal_2013}.
\contribution{This is the first chapter of my student's (Ms. Helm) PhD dissertation. The work is funded with a Brown/MBL seed grant that I secured with J Smith. I helped design the project, implement the analyses, and write the manuscript. Most of the work for this project was done in my lab.}

\iftoggle{long}{
\item Howison, M, N Sinnott-Armstrong, CW Dunn (2012) BioLite, a lightweight bioinformatics framework with automated tracking of diagnostics and provenance. 4th USENIX Workshop on the Theory and Practice of Provenance (TaPP12). \url{https://www.usenix.org/system/files/conference/tapp12/tapp12-final5.pdf} (peer-reviewed conference proceeding). \contribution{I conceived of this software package and secured support for its development. I helped write the code and plan the architecture of the tool, as well as write the manuscript. The other two authors are members of my lab.}
}

\item Smith, SA, NG Wilson, F Goetz, C Feehery, SCS Andrade, GW Rouse, G Giribet, CW Dunn (2011) Resolving the evolutionary relationships of molluscs with phylogenomic tools. Nature 480:364-367. \url{http://dx.doi.org/10.1038/nature10526}. \contribution{This study was largely conducted in my lab. I conceived of the project with NGW and GG. I designed and oversaw much of the data acquisition strategy. I designed the data analysis strategy with SAS. I wrote the majority of the manuscript.}

\item Siebert, S,  MD Robinson, SC Tintori, F Goetz, RR Helm, SA Smith, N Shaner, SHD Haddock, CW Dunn (2011) Differential Gene Expression in the Siphonophore \emph{Nanomia bijuga} (Cnidaria) Assessed with Multiple Next-Generation Sequencing Workflows. PLoS One 6(7): e22953. \url{http://dx.doi.org/10.1371/journal.pone.0022953}. \contribution{This study was largely conducted in my lab. I conceived of the project design, much of the data acquisition strategy, and collaborated on the data analysis (including writing new software). I wrote much of the manuscript.}

\iftoggle{long}{
\item Edgecombe, GD, G Giribet, CW Dunn, A Hejnol, RM Kristensen, RC Neves, GW Rouse, K Worsaae, and MV S{\o}rensen (2011) Higher-level metazoan relationships: recent progress and remaining questions. Organisms, Diversity, and Evolution 11:151-172. \url{http://dx.doi.org/10.1007/s13127-011-0044-4}. \contribution{I conducted the dating analyses, made figures 2 and 4, wrote the first draft of several sections, and helped revise the manuscript. }
}

\item Hejnol, A, M Obst, A Stamatakis, M Ott, G Rouse, G Edgecombe, P Martinez, J Bagu\~{n}\`{a}, X Bailly, U Jondelius, M Wiens, WEG M\"{u}ller, Elaine Seaver, WC Wheeler, MQ Martindale, G Giribet, and CW Dunn (2009) Assessing the root of bilaterian animals with scalable phylogenomic methods. Proc. R. Soc. B. 276:4261-4270. \url{http://dx.doi.org/10.1098/rspb.2009.0896}. Git repository: \url{https://bitbucket.org/caseywdunn/hejnol_etal_2009}. \contribution{I designed much of the project, developed novel analytic methods, did most of the analyses, produced figures, and wrote much of the manuscript. }

\iftoggle{long}{
\item Cartwright, P, NM Evans, CW Dunn, AC Marques, MP Miglietta, P Schuchert, and AG Collins (2008) Phylogenetics of Hydroidolina (Hydrozoa: Cnidaria). Journal of the Marine Biological Association of the United Kingdom 88:1663-1672. \url{http://dx.doi.org/10.1017/S0025315408002257}. \contribution{I contributed siphonophore data and helped to revise the manuscript. }
}

\item Dunn, CW, A Hejnol, DQ Matus, K Pang, WE Browne, SA Smith, E Seaver, GW Rouse, M Obst, GD Edgecombe, MV S{\o}rensen, SHD Haddock, A Schmidt-Rhaesa, A Okusu, RM Kristensen, WC Wheeler, MQ Martindale, and G Giribet (2008) Broad phylogenomic sampling improves resolution of the Animal Tree of Life. Nature 452:745-749. \url{http://dx.doi.org/10.1038/nature06614}. (Cover Article). \contribution{I prepared many samples for sequencing, developed novel analytic methods, wrote extensive software tools to implement these new methods, did most of the analyses, produced figures, and wrote most of the manuscript. Though this project was initiated during my postdoc, much of it was done after arriving at Brown. }

\iftoggle{long}{
\item Smith, SA and CW Dunn (2008) Phyutility: a phyloinformatics tool for trees, alignments, and molecular data. Bioinformatics. 24:715-716. \url{http://dx.doi.org/10.1093/bioinformatics/btm619}. Code repository: \url{https://code.google.com/p/phyutility} . \contribution{I helped conceive of this software tool, assisted in design decisions, and helped to write the manuscript.}


\item Giribet, G, CW Dunn, GD Edgecombe, and GW Rouse (2007) A modern look at the Animal Tree of Life. Zootaxa 1668:61-79. \contribution{I contributed to several sections of the paper, helped with figures, and helped revise the manuscript. }

\item Oota, H, CW Dunn, WC Speed, AJ Pakstis, MA Palmatier, JR Kidd and KK Kidd (2007) Conservative evolution in duplicated genes of the primate Class I ADH cluster. Gene 392:64-76. \url{http://dx.doi.org/10.1016/j.gene.2006.11.008}. \contribution{I performed several phylogenetic analyses. }
}

\item Dunn, CW and GP Wagner (2006) The evolution of colony-level development in the Siphonophora (Cnidaria:Hydrozoa). Development, Genes, and Evolution. 216:743-754. \url{http://dx.doi.org/10.1007/s00427-006-0101-8} (Cover Article). \contribution{ This is a chapter of my PhD thesis. I collected all the data and wrote the manuscript. }

\item Matus, DQ, K Pang, H Marlow, CW Dunn, GH Thomsen, and MQ Martindale (2006) Deep evolutionary roots for bilaterality in the metazoa. Proceedings of the National Academy of Sciences USA. 103:11195-11200. \url{http://dx.doi.org/10.1073/pnas.0601257103}. \contribution{ I performed several bioinformatics analyses and helped to write portions of the manuscript. }

\item Matus, DQ, RR Copley, CW Dunn, A Hejnol, H Eccleston, KM Halanych, MQ Martindale, and MJ Telford (2006) Broad Taxon and Gene Sampling Indicate that Chaetognaths Are Protostomes. Current Biology. 16:R575-R576. \url{http://dx.doi.org/10.1016/j.cub.2006.07.017}. \contribution{ I performed several bioinformatics analyses and helped to write portions of the manuscript. }

\iftoggle{long}{
\item Dunn, CW, PR Pugh, and SHD Haddock (2005) Molecular phylogenetics of the Siphonophora (Cnidaria), with implications for the evolution of functional specialization. Systematic Biology 54:916-935. \url{http://dx.doi.org/10.1080/10635150500354837} (Cover Article). Git repository: \url{https://bitbucket.org/caseywdunn/siphonophores_2005}. \contribution{This is a chapter from my PhD thesis. I designed the project, collected many of the specimens, collected all the DNA sequence data, analyzed the data, and wrote the manuscript. }

\item Dunn, CW (2005) The complex colony-level organization of the deep-sea siphonophore \emph{Bargmannia elongata} (Cnidaria, Hydrozoa) is directionally asymmetric and arises by the subdivision of pro-buds. Developmental Dynamics 234:835-845. \url{http://dx.doi.org/10.1002/dvdy.20483} (Cover Article). \contribution{This is a chapter from my PhD thesis. I was the sole contributor to the study.}
}

\item Haddock, SHD, CW Dunn, PR Pugh and CE Schnitzler (2005) Bioluminescent and red-fluorescent lures in a deep-sea siphonophore. Science 309:263. \url{http://dx.doi.org/10.1126/science.1110441}. \contribution{I made the initial observation of the lures, helped describe their morphology, and participated in writing the manuscript.}

\iftoggle{long}{
\item Dunn, CW, PR Pugh, and SHD Haddock (2005) \emph{Marrus claudanielis}, a new species of deep-sea physonect siphonophore (Siphonophora, Physonectae). Bulletin of Marine Science 76:699-714. \contribution{This is a chapter from my PhD thesis. I made many of the observations and illustrations, and took the lead writing the manuscript. }

\item Haddock, SHD, CW Dunn, and PR Pugh (2005) A reexamination of siphonophore terminology and morphology, applied to the description of two new prayine species with remarkable bio-optical properties. Journal of the  Marine Biological Association of the  U.K. 85:695-707. \url{http://dx.doi.org/10.1017/S0025315405011616}. \contribution{This is a chapter from my PhD thesis. I contributed to the observations, illustrations, and writing. I took the lead in developing the new morphological terminology. }

\item Lynch, VJ, JJ Roth, K Takahashi, CW Dunn, DF Nonaka, G Stopper and GP Wagner (2004) Adaptive evolution of HoxA-11 and HoxA-13 at the origin of the uterus in mammals. Proceedings of the Royal Society B: Biological Sciences 271:2201-2207. \url{http://dx.doi.org/10.1098/rspb.2004.2848}. \contribution{I collected data from armadillo.}
}

\end{itemize}

\iftoggle{long}{


\subsection*{d. Non-refereed journal articles}

\begin{itemize}

\item Lewis, ZR and CW Dunn. (2018) Genome evolution: We are not so special. eLife. \url{http://dx.doi.org/10.7554/eLife.38726}.  Git code repository: \url{https://github.com/dunnlab/gene_inventory_2018}.

\item Dunn, CW (2017) Ctenophore trees. Nature Ecology and Evolution. \url{http://dx.doi.org/10.1038/s41559-017-0359-4}

\item Hejnol, A, CW Dunn (2016) Animal Evolution: Are Phyla Real? Current Biology 26:R424-R426. \url{http://dx.doi.org/10.1016/j.cub.2016.03.058}.

\item Dunn, CW (2014) Keeping Mates Close and Competition Out in an Ocean Sponge, November 13, 2014. \url{http://www.nytimes.com/2014/11/13/science/keeping-mates-close-and-competition-out-in-an-ocean-sponge.html}

\item Dunn, CW (2014) Fiery Bodies Under the Waves. New York Times, August 14, 2014. \url{http://www.nytimes.com/2014/08/14/science/a-colonizing-fire-of-giant-plankton.html}

\item Dunn, CW (2014) A Marine Magician's Vanishing Act. New York Times, May 28, 2014. \url{http://www.nytimes.com/2014/05/28/science/a-marine-magicians-vanishing-act.html}

\item Dunn, CW (2014) A Cloak of Near Invisibility in an Underwater World. New York Times, April 24, 2014. \url{http://www.nytimes.com/2014/04/24/science/a-cloak-of-near-invisibility-in-an-underwater-world.html}

\item Dunn, CW (2014) Two Urchins, Similar but Not. New York Times, February 26, 2014. \url{http://www.nytimes.com/2014/02/26/science/two-urchins-similar-but-not.html}

\item Dunn, CW (2014) Poisonous Prey Turned Into Hunter's Defense. New York Times, February 12, 2014. \url{http://www.nytimes.com/2014/02/13/science/poisonous-prey-turned-into-hunters-defense.html}

\item Dunn, CW (2014) Moving, Without Feet to Do So. New York Times, January 23, 2014. \url{http://www.nytimes.com/2014/01/23/science/earth/moving-without-feet-to-do-so.html}

\item Dunn, CW (2013) The Color of Royalty, Bestowed by Science and Snails. New York Times, October 9, 2013. \url{http://www.nytimes.com/2013/10/09/science/the-color-of-royalty-bestowed-by-science-and-snails.html}

\item Dunn, CW (2013) As 'Normal' as Rabbits' Weights and Dragons' Wings. New York Times, September 23, 2013. \url{http://www.nytimes.com/2013/09/24/science/as-normal-as-rabbits-weights-and-dragons-wings.html}

\item Dunn, CW (2013) Sex in Spoonworms. New York Times, September 16, 2013. \url{http://www.nytimes.com/2013/09/17/science/creatures-strange-and-complex-in-colorful-detail.html} (the text I wrote accompanies the video).

\item Dunn, CW (2013) Evolution: Out of the ocean. Current Biology 23:R242-R243. \url{http://dx.doi.org/10.1016/j.cub.2013.01.067}. \contribution{I was invited to submit this dispatch article to summarize our current understanding of when animals colonized land.}

\item Dunn, CW (2009) Siphonophores. Current Biology 19:R233-R234. \url{http://dx.doi.org/10.1016/j.cub.2009.02.009}. \contribution{I was invited to submit this review, and was the sole author.}

\item Haddock, SHD and CW Dunn (2005) The complex world of siphonophores. JMBA Global Marine Environment 2005(2):24-25.  \contribution{I helped write this perspectives piece.}

\end{itemize}


\subsection*{f. Abstracts}

\begin{itemize}

\item Dunn, CW (2013) The comparative biology of gene expression. Invited speaker for the symposium ``Understanding First Order Phenotypes: Transcriptomics for Emerging Models'' at the 2013 meeting for the Society of Integrative and Comparative Biology. Abstract S4-2.1. \url{http://www.sicb.org/meetings/2013/SICB\%202013\%20abstracts.pdf}. \contribution{This talk presented theoretical work on the evolutionary analysis of large functional genomic datasets. I developed most of the work, and prepared and presented the talk.}

\item Helm, RR and CW Dunn (2013) The evolution of direct development in Scyphozoa. A poster presented at the 2013 meeting for the Society of Integrative and Comparative Biology. Abstract P1.58. \url{http://www.sicb.org/meetings/2013/SICB\%202013\%20abstracts.pdf}. \contribution{This poster was prepared and presented by RR Helm, a graduate student in my lab. It presented thesis work she has conducted in my lab in collaboration with me.}


\end{itemize}



\subsection*{g. Invited lectures}

\begin{entry}{60pt}

\item [\rm 2023] Biodiversity Cell Atlas Meeting (Barcelona, Spain), International Congress on Invertebrate Morphology (Vienna, Austria)

\item [\rm 2022] Moss Landing Marine Lab, Jacques Monod Conference Metazoan Origins (Roscoff, France)

\item [\rm 2019] Keynote speaker at the Simon Fraser University Day of Omics, Systems Biology Teory Lunch invited speaker (Harvard), Graduate Student Invited Speaker (Stony Brook, Department of Ecology and Evolution), Yale University Department of Genetics

\item [\rm 2018] UC Davis Department of Evolution and Ecology, University of Guelph Department of Integrative Biology

\item [\rm 2017] Harvard Department of Organismal and Evolutionary Biology graduate student invited speaker, University of Alabama Allele Series speaker, Missouri Botanical Garden Annual Fall Symposium invited speaker, University of Texas

\item [\rm 2016] Stanford University Abbott Lecture, Ctenopalooza (University of Florida) Keynote Speaker, University of Maryland BISI-BEES seminar student invited speaker, Blavatnik Science Symposium (New York Academy of Sciences), Yale University Department of Ecology and Evolutionary Biology, Boston Evolutionary Genomics Supergroup meeting at Harvard

\item [\rm 2015] Gordon Research Conference, Friday Harbor Marine Laboratory, Systematics and Biodiversity Symposium at the Norwegian Academy of Sciences and Letters, University of Rhode Island

\item [\rm 2014] Universidad Austral de Chile (Valdivia, Chile), Bowdoin College (Brunswick, Maine), University of Missouri, Evolution of First Nervous Systems II meeting (Whitney Laboratory for Marine Bioscience, Florida), 13th Annual Genome Symposium (New York University), Kyushu University (Fukuoka, Japan)

\item [\rm 2013] American Museum of Natural History, Global Invertebrate Genomics Alliance (GIGA) workshop at Nova University (Fort Lauderdale, FL), meeting of the Centre National de Ressources Biologiques Marines, EMBRC (Villefranche, France), Brown Department of Ecology and Evolutionary Biology Seminar, Brown Computer Science Department Industrial Partners Program Symposium, Tree of Life Symposium at the 2013 meeting of the Society of Systematic Biologists (title: ``Building trees with transcriptomes and analyzing transcriptomes with trees''), Brown EPSCoR Bioinformatics Workshop

\item [\rm 2012] World Economic Forum Annual Meeting (Davos, Switzerland), Smithsonian Tropical Research Institute (Panama), Sars International Centre for Marine Molecular Biology (Bergen, Norway), Clark University, University of Miami, Stony Brook University, Applied Math at Brown University, University of Florida

\item [\rm 2011] National Science Board, (Washington, DC), National Science Foundation (Washington, DC), Marine Biological Laboratory (Woods Hole, MA), Boston University (Boston, MA), Marine Observatory at Villefranche-sur-Mer (France), Smithsonian National Museum of Natural History Department of Invertebrate Zoology (Washington, DC), Smithsonian National Museum of Natural History - Symposium ``Next generation Sequencing: Transformative Technology for Biodiversity Science'', keynote speaker at the Brown University Communicating Your Science Workshop, Rhode Island College

\item [\rm 2010] Harvard University, University of Gothenburg (Sweden), Iowa State University, Brown University Wayland Collegium, Helicos BioSciences, National Academy of Science ``Kavli Frontiers of Science'' meeting, PopTech, University of Texas at Austin

\item [\rm 2009] Yale University, University at Buffalo, LM University Munich, meeting of the German Research Foundation Priority Program ``Deep Metazoan Phylogeny'' at Humboldt University Berlin, University of Connecticut, Harvard University course OEB275r: ``Phylogenetics in the Era of Genomics'', University of Rhode Island, Barcelona University, Brown MCB, invited speaker at the Society for Systematic Biologists symposium ``Advances in Tree Reconstruction from Complex Data Matrices'' in Moscow, Idaho

\item [\rm 2008] University of Rochester, Friday Harbor Marine Laboratory, Woods Hole Oceanographic Institute, New England Biolabs

\item [\rm 2007] Harvard University, Roger Williams University

\item [\rm 2005] University of Hawaii

\end{entry}



\subsection*{i. Work in review}

This section includes publicly available preprints.

\subsection*{j. Patents and inventions}

\begin{entry}{60pt}

\item [\rm 2016] US Patent number 9,330,295. ``Spatial sequencing/gene expression camera''

\end{entry}

\iftoggle{inreview}{



}% end \iftoggle{inreview}







\iftoggle{inprep}{

\subsection*{i. Work in progress}

This list includes selected manuscripts that are close to submission. The author lists and journals may change prior to publication.


}% end \iftoggle{inprep}

} % end \iftoggle{long}


\iftoggle{grants}{

\iftoggle{long}{
\section{Research Grants}
}{
\section{Selected Research Grants}
} %\iftoggle{long}

\subsection*{a. Current grants}

\begin{itemize}

\item YIBS Biodiversity and Ecosystems Award, 2022, \$180,000

\end{itemize}

\subsection*{b. Completed grants}

\begin{itemize}

\iftoggle{long}{

\item ``Collaborative research: The effects of predator traits on the structure of oceanic food webs'', NSF AWD0001810, September 1, 2018 - August 31, 2021, \$350,533.00, Principal Investigator.

\item ``Real Time Phylogeny and Contact Tracing to Disrupt HIV Transmission'', NIH R01AI136058 subcontract via Miriam Hospital, 2018-2023. \$214,801

\item ``Ginsberg Award, Peabody Museum, 2018-2021''. \$18,000

\item Alan T. Waterman Award, NSF, May, 2011. \$500,000 (no indirect charges)

\item ``Inferring HIV Transmission Networks from Deep Viral Phylogenies'', Center for AIDS Research, January 1, 2017 - December 31, 2017, \$40,000, Principal Investigator.

\item ``The evolution of gene expression and functional specialization in Siphonophora (Cnidaria)'', NSF, March 2013- March 2017. \$700,000, Principal Investigator.

\item ``Extended phylogenetic Inference of HIV Transmission Network'', Brown University BioMed Emerging Areas of New Science Deans Awards, 2014-2015. \$80,000, co-Principal Investigator with Rami Kantor.

\item ``Making sense of the data windfall: New statistical approaches to evolutionary analyses of gene expression'', Brown University Office of the Vice President of Research, 2013-2014. \$100,000, co-Principal Investigator with Xi Luo and Zhijin Wu.
} %\iftoggle{long}

\item ``Collaborative Research: Resolving old questions in Mollusc phylogenetics with new EST data and developing general phylogenomic tools'', NSF, submitted July 2008, February 2009- January 2013. \$450,000 (\$775,000 total), Principal Investigator.

\iftoggle{long}{
\item ``IGERT:  Reverse Ecology: Computational Integration of Genomes, Organisms and Environments'', NSF, submitted September, 2009. Participant.

\item Brown/MBL Seed Fund Award (with Joel Smith). \$12,500. Funded August, 2010.

\item ``Phylogenetic analyses of quantitative differential expression data'' Brown CCMB Seed fund. \$5,000. Funded October, 2010.
}  %\iftoggle{long}

\item ``PSCIC Full Proposal: The iPlant Collaborative: A Cyberinfrastructure-Centered Community for a New Plant Biology'', NSF, subcontract via University of Arizona, July 1, 2009 -- June 30, 2012. \$257,762.

\iftoggle{long}{
\item ``DISSERTATION RESEARCH: Phylogeny and the Evolution of Colony Organization in the Siphonophora, Cnidaria'', NSF, June 1, 2004 -- December 31, 2005, \$11,985, Co-Principal Investigator.
} %\iftoggle{long}

\end{itemize}


\iftoggle{declinedgrants}{

\subsection*{c. Proposals submitted (other than above)}

\begin{itemize}

\item ``Collaborative Research: Phylogenetic approaches to comparative functional genomics'', NSF, submitted August 2017 (declined). Principal Investigator.
\item ``COLLABORATIVE RESEARCH: ABI Innovation: Statistical measures of uncertainty and heterogeneity for phylogenetic analysis of genomic data'', NSF, submitted September 2017 (pending). co-Principal Investigator.
\item ``Preliminary Proposal: Collaborative Research: RUI: An integrated phylogenomic analysis of animal gene family evolution'', NSF, submitted January 2016 (declined). co-Principal Investigator.
\item ``Blavatnik - Awards for Young Scientists'', nominated November 2015 (declined)
\item ``HHMI - Faculty Scholars Competition'', July 2015 (declined)
\item ``Data-Driven Discovery Initiative'', Betty and Gordon Moore Foundation preproposal, submitted March 2014 (my preproposal was one of the 93 that was selected for a full proposal from the 1095 submitted; my full proposal declined). Principal Investigator.
\item ``ABI Innovation: The theory, implementation, and application of Bayesian transcriptome assembly'', NSF, submitted August 2014 (declined). Principal Investigator.
\item ``IOS Preliminary proposal: The colony-level development of Siphonophora (Cnidaria)'', NSF, submitted January 2014 (declined). Principal Investigator.
\item ``IOS Preliminary Proposal: Collaborative Research: Identifying the genomic and transcriptomic basis for bioluminescence and light detection in ctenophores'', NSF, submitted January 2014 (declined). co-Principal Investigator.
\item ``ABI Innovation: Bayesian genome assembly and assessment by Markov Chain Monte Carlo sampling'', NSF, submitted August 2013 (declined). Principal Investigator.
\item ``DISSERTATION RESEARCH: The evolution of complex lifecycles in Scyphozoa'', NSF, submitted November 2012 (declined). Principal Investigator.
\item ``Collaborative Proposal: Assembling the Animal Tree of Life'', NSF, submitted March, 2012 (declined). \$1,301,662. Principal Investigator.
\item ``The evolution of gene expression and functional specialization in siphonophora (Cnidaria)'', NSF, submitted January 10, 2011 (declined). \$777,916.
\item ``Packard fellowship'', Packard Foundation, submitted April 8, 2010 (declined). \$875,000. Principal Investigator.
\item W. M. Keck Foundation Medical Research Program Preproposal, submitted April, 2010 (declined). \$993,339. Principal Investigator.
\item ``Untangling the Genotype-Phenotype Map With Phylogenetic Investigations of Gene Expression'', Pew Charitable Trusts, submitted November 2, 2009 (declined). \$240,000. Principal Investigator.
\item ``Collaborative Research: Emergence of the Germ Line from Mesoderm: A Transcriptome Approach to Cell Type Specification and Evolution'', NSF, submitted June 13, 2009 (declined), \$382,889. Principal Investigator.
\item ``Packard fellowship'', Packard Foundation, submitted April 20, 2009 (declined). \$875,000. Principal Investigator.
\item ``Collaborative Research: Assembling the Animal Tree of Life: Resolving the root and key nodes'', NSF, submitted March 20, 2009 (declined). \$600,574. Principal Investigator.
\item ``Collaborative Research: Resolving old questions in Mollusc phylogenetics with new EST data and developing general phylogenomic tools'', NSF, submitted January 7, 2008 (declined). \$616,170. Principal Investigator.
\item ``PEET: Monographic studies and taxonomic training on Siphonophora (Cnidaria)''. NSF, submitted March 5, 2007 (declined). \$750,000. Principal Investigator.



\end{itemize}

} % end \iftoggle{declinedgrants}

} % end \iftoggle{grants}




\iftoggle{service}{

\section{Service}

\subsection{To the University}

\begin{entry}{60pt}

\item [\rm 2024] EEB Director of Graduate Studies (DGS), Humanities Advisory and Tenure and Appointments Committee (HTAC)

\item [\rm 2023] EEB Director of Graduate Studies (DGS), CEID review committee

\item [\rm 2022] EEB Director of Graduate Studies (DGS), Committee on the Economic Status of the Faculty (CESOF), EEB safety officer (overseeing COVID research reactivation and approval), Yale Genetics Training Grant (TPG) Executive Committee, Undergraduate Admissions, Yale Institute for Biospheric Studies (YIBS) small grant director

\item [\rm 2021] EEB Director of Graduate Studies (DGS), Committee on the Economic Status of the Faculty (CESOF), EEB safety officer (overseeing COVID research reactivation and approval), Yale Genetics Training Grant (TPG) Executive Committee, Yale Institute for Biospheric Studies (YIBS) small grant director

\item [\rm 2020] Biological Sciences Advisory Committee (July 1, 2018, to June 30, 2020), Life Sciences department review at Yale NUS, faculty search committee at Yale NUS, Committee on the Economic Status of the Faculty (CESOF), EEB safety officer (overseeing COVID research reactivation and approval), Yale Institute for Biospheric Studies (YIBS) small grant director

\item [\rm 2019] Biological Sciences Advisory Committee (July 1, 2018, to June 30, 2020), Multiple committees at Yale NUS

\item [\rm 2018] Biological Sciences Advisory Committee (July 1, 2018, to June 30, 2020), Peabody Exhibits Committee

The following were at Brown University:

\item [\rm 2016] Served on the Digital Teaching and Learning Committee, Director of the COBRE Computational Biology Core

\item [\rm 2014] Served on the Steering Committee for the Brown Genomics Core Facility, served on the Advisory Board for the Brown Center for Computation and Visualization, maintained department website for EEB, served on the the provost's committee to explore RISD/Brown collaboration, guest lecture in BIOL2010

\item [\rm 2013] Served on the University Strategic Planning Committee on Online Teaching and Learning, served on the Steering Committee for the Brown Genomics Core Facility, served on the Advisory Board for the Brown Center for Computation and Visualization, assisted in the creation of a new department website for EEB, served on the the provost's committee to explore RISD/Brown collaboration

\item [\rm 2012] Participated in the Strategic Planning Sessions for the Program in Biology, served on the University Strategic Planning Committee on Online Teaching and Learning, served on the Steering Committee for the Brown Genomics Core Facility, served on the Advisory Board for the Brown Center for Computation and Visualization, presented at the Academic Technology Showcase, filed New Invention Disclosures (Tech ID 2157, 2174)

\item [\rm 2011] Co-directed the Center for Computational Molecular Biology ``Genome Assembly Special Forces Workshop'', member of the CCV Advisory Board for High Performance Computing, member of the Illumina steering committee, was a presenter to the visiting World Economic Forum panel, presented to the Medical School Advisory Council

\item [\rm 2010] Directly involved in hiring a bioinformatician in BioMed, faculty representative on panel to redesign the Brown home page, played an active role (including technical support in the lab) on the selection of a next-generation sequencer for the NIH SIG award, member of the CCV Advisory Board for High Performance Computing, member of the Illumina steering committee, participated in the EEB internal departmental review. Teaching and mentorship: sophomore advisor to two undergraduate students, MCB trainer.

\item [\rm 2009] Participated in the OVPR sponsored ``Brown Research Data Storage Focus Group'', provided assistance with NIH SIG proposal for a next-generation DNA sequencer, participated in discussions with IBM regarding High Performance Computing acquisition, panel member at the Sheridan Center roundtable ``UTRAs -- Integrating Research and Teaching in the Life \& Physical Sciences'', joined the CCV Advisory Board for High Performance Computing at the invitation of the Vice President for Research, played an active role in next-generation sequencing platform selection for NIH SIG Award. Teaching and mentorship: was a sophomore advisor to two students, became a graduate student trainer in MCB, organized a special campus-wide presentation by the science writer Carl Zimmer sponsored by EEB, BIBS, and ECI

\item [\rm 2008] Assisted in the preparation of a Rhode Island EPSCoR proposal for next generation sequencing and downstream computational analyses; Participated in discussions with IBM to develop supercomputing resources at Brown


\end{entry}



\subsection{To the profession}

Editorial
\begin{entry}{60pt}
\item [\rm 2025] Guest editor PLoS Computational Biology
\item [\rm 2015] Guest editor at Proceedings of the National Academy of Science USA (PNAS)
\item [\rm 2014] Guest editor at Proceedings of the National Academy of Science USA (PNAS), co-hosted (with Chris Pires) the Society for Systematic Biologists symposium ``Phylogenomics, transcriptomics, and the evolution of gene expression'' at the annual Evolution meeting
\item [\rm 2012-present] Member of the editorial board of the Journal of Experimental Zoology Part B: Molecular and Developmental Evolution
\end{entry}

Journal Reviews
\begin{entry}{60pt}
\item [\rm 2025] Current Biology, Nature, Nature Communications, PNAS, Mitochondrial DNA: Resources, Molecular Biology and Evolution, Current Opinion in Genetics and Development
\item [\rm 2024] Molecular Biology and Evolution, Current Biology, Proceedings of the Royal Academy Series B, PNAS
\item [\rm 2023] Bioinformatics Advances, Molecular Biology and Evolution, Genome Research, Evolution and Development
\item [\rm 2022] PNAS, Nature Ecology and Evolutionary Biology, Nature Communications, BMC Biology, Bioinformatics, Science, Genome Biology
\item [\rm 2021] PeerJ, ICOES, MBE, Nature
\item [\rm 2020] PNAS, Science
\item [\rm 2019] Nature, Developmental Biology, Development, PNAS
\item [\rm 2018] Systematic Biology, Nature, Current Biology, Nature Ecology and Evolutionary Biology, Molecular Biology and Evolution, PeerJ, Organisms Development and Evolution
\item [\rm 2017] Systematic Biology, Nature, Nature Ecology and Evolution, Nature Scientific Reports, BMC Evolutionary Biology
\item [\rm 2016] Systematic Biology, Nature Scientific Reports, Current Biology, Journal of Experimental Zoology Series B, BMC Evolutionary Biology, eLIFE, PNAS
\item [\rm 2015] PNAS, Current Biology, Nature Communications, EvoDevo, Systematic Biology, Proceedings of the Royal Society Series B, Nature Protocols, Nature, Zootaxa
\item [\rm 2014] PLoS Currents, PNAS
\item [\rm 2013] BMC Genomics, Nature, Genome Biology and Evolution, Current Biology, PLoS One, Molecular Biology and Evolution
\item [\rm 2012] Nature, PNAS, Bioinformatics, Current Biology, Molecular Reproduction and Development
\item [\rm 2011] Current Biology, Molecular Biology and Evolution, Evolutionary Ecology, Mitochondrial DNA, PLoS One
\item [\rm 2010] Proceedings of the National Academy of Sciences (USA), Current Biology, BMC Evolutionary Biology, BMC Bioinformatics, PLoS Biology, Molecular Phylogenetics and Evolution, Trends in Genetics, Biology Letters
\item [\rm 2009] Nature, Molecular Biology and Evolution, Current Biology, Zootaxa, BMC Evolutionary Biology, Molecular Phylogenetics and Evolution
\item [\rm 2008] Current Biology, BMC Evolutionary Biology, Systematic Biology, Proceedings of the Royal Society B, Journal of the Marine Biological Association of the United Kingdom, Caribbean Journal of Science, Deep-Sea Research Part I, Journal of Molecular Biology, Bioinformatics
\item [\rm 2007] Developmental Dynamics, Hydrobiologia, Integrative and Comparative Biology, JEZ Part B: Molecular and Developmental Evolution, Journal of the Marine Biological Association of the UK, Systematics and Biodiversity, Zoologica Scripta, Zootaxa

\end{entry}
Grant Reviews and Funding Agency Activities
\begin{entry}{60pt}
\item [\rm 2021] Reviewed proposals for NSF
\item [\rm 2016] Served on the NSF Evolution of Developmental Mechanisms grant review panel
\item [\rm 2015] Reviewed proposals for NSF
\item [\rm 2014] Served on the NSF Phylogenetic Systematics grant review panel
\item [\rm 2013] German-Israeli Foundation for Scientific Research and Development (GIF), NSF
\item [\rm 2012] Reviewed proposals for NSF, Israel Science Foundation, Sloan Foundation
\item [\rm 2011] Reviewed multiple NSF grant proposals
\item [\rm 2010] Reviewed multiple NSF grant proposals, reviewed grant proposal for the United Kingdom Biotechnology And Biological Sciences Research Council
\item [\rm 2009] Grant review panelist for the NSF Systematic Biology and Biological Inventories Cluster, participant in the NSF workshop ``Future Directions in Biodiversity and Systematics Research''
\item [\rm 2008] Reviewed an NSF grant proposal, assisted the NSF with a special report on jellyfish


\end{entry}
Other Activities
\begin{entry}{60pt}
\item [\rm 2019- present] Trustee of the Society of Systematic Biologists
\item [\rm January 1, 2016- December 31, 2018] Treasurer of the Society of Systematic Biologists
\item [\rm 2013-2916] Member of the supervisory board of the Neptune Consortium (\url{http://neptune-itn.eu/}), a multi-institution evolutionary developmental biology and neurobiology training network funded by the European Union
\item [\rm 2012] Served as a faculty member for the PopTech Science Fellows Program

\end{entry}


\subsection{To the community}

\begin{entry}{60pt}

\item [\rm 2012] Participation in the World Economic Forum Meeting at Davos
\item [\rm 2011] Furthered development of creaturecast.org, hosted a RISD class in my lab
\item [\rm 2010] Furthered development of creaturecast.org, which continues to have wide distribution. Videos have now been viewed over 200,000 times, the site is cross posted by nature.com. Interviewed or had work featured in magazines for a public audience, including Brown Alumni Magazine, Brown Medicine Magazine, The Scientist
\item [\rm 2009] Launched creaturecast.org, a collaborative blog and podcast series aimed at communicating biological research to a broad audience. It has been featured on National Public Radio, iTunes, Voice of America, and at a wide variety of other outlets. Videos produced as part of this project have been viewed more than 100,000 times.
\item [\rm 2008] Reviewed portions of the introductory Biology textbook ``Life: The Science of Biology, 8th Edition'' (Sinauer and Associates) for the editor
\item [\rm Previous] I wrote and maintain siphonophores.org, a web site that provides an introduction to siphonophore biology for general and scientific audiences. Siphonophores.org was featured in the August 12, 2005 NetWatch section of Science.

\end{entry}


\subsection{Selected media, reviews of work, and press engagement}

\begin{entry}{60pt}

\item [\rm 2017] Interviewed for \emph{``Ecologists protest sudden end of NSF dissertation grants''. Science, June, 2017.}(\url{http://www.sciencemag.org/news/2017/06/ecologists-protest-sudden-end-nsf-dissertation-grants})
\item [\rm 2017] Interviewed for \emph{``Big data renews fight over animal origins''. Nature, March 24, 2017.}(\url{http://www.nature.com/news/big-data-renews-fight-over-animal-origins-1.21703})
\item [\rm 2015] Interviewed for \emph{``Is That a Colony of Conjoined Twins or a Single Sea Creature?''. The Atlantic, September 4, 2015.}(\url{http://www.theatlantic.com/technology/archive/2015/09/living-jet-engines-colony-nanomia-siphonophore/403816/})
\item [\rm 2015] Article that features our paper on fluorescent lures. \emph{``Light show lures prey''. Nature, August 12, 2015.}(\url{http://www.nature.com/nature/journal/v524/n7564/full/524138a.html})
\item [\rm 2015] Article that features our paper on fluorescent lures. \emph{Randall, Ian. ``Slideshow: Glowing predators of the deep''. Science, August 7, 2015.}(\url{http://dx.doi.org/10.1126/science.aac8986})
\item [\rm 2015] Article that features our analysis of siphonophore data using the Amazon cloud. \emph{Drake, Nadia. ``How to catch a cloud''. Nature, June 3, 2015.}(\url{http://www.nature.com/news/how-to-catch-a-cloud-1.17672})
\item [\rm 2015] Article about our paper on sponges and ctenophores. \emph{Yong, Ed. ``Consider the Sponge''. New Yorker Elements, April 24, 2015.}(\url{http://www.newyorker.com/tech/elements/consider-the-sponge})
\item [\rm 2014] Article about a new group of animals. \emph{Skinner, Nicole. ``Sea creatures add branch to tree of life''. Nature, September 3, 2014.}(\url{http://dx.doi.org/10.1038/nature.2014.15833})
\item [\rm 2014] Interviewed for \emph{Zimmer, Carl. ``Strange Findings on Comb Jellies Uproot Animal Family Tree''. National Geographic, May 21, 2014.} (\url{http://news.nationalgeographic.com/news/2014/05/140521-comb-jelly-ctenophores-oldest-animal-family-tree-science/})
\item [\rm 2014] Article about ctenophore research and the evolution of complexity.  \emph{Maxmen, Amy. ``Evolution, You're Drunk''. Nautilus, January 30, 2014.}  (\url{http://nautil.us/issue/9/time/evolution-youre-drunk})
\item [\rm 2014] Article about ctenophore research. \emph{Nadboy, Jason. ``Humans most distant animal relative found''. Brown Daily Herald, January 29, 2014.} (\url{http://www.browndailyherald.com/2014/01/29/humans-distant-animal-relative-found/})
\item [\rm 2013] Interviewed on Rhode Island Public Radio about ctenophore research. (\url{http://ripr.org/post/local-comb-jelly-our-closest-ancient-relative})
\item [\rm 2013] Interviewed for \emph{Barber, Elizabeth. ``Is your oldest sister a warty comb jelly?''. The Christian Science Monitor, December 12, 2013.}  (\url{http://www.csmonitor.com/Science/2013/1212/Is-your-oldest-sister-a-warty-comb-jelly})
\item [\rm 2013] Interviewed for \emph{Karberg, Sascha. ``Evolution der Tiere: Am Anfang war der Glibber''. Der Spiegel, December 13, 2013.}  (\url{http://www.spiegel.de/wissenschaft/natur/rippenquallen-genom-entschluesselt-hinweise-auf-stammbaum-der-tiere-a-938661.html})
\item [\rm 2013] Article about CreatureCast.  \emph{``Animals Get the Spotlight in CreatureCast''. Science, October 18, 2013.} (\url{http://www.sciencemag.org/content/342/6156/294.3.full})
\item [\rm 2013] Article about CreatureCast. \emph{Friedmann, Meghan. ``CreatureCast videos featured on NY Times website''. Brown Daily Herald, October 2, 2013.} (\url{http://www.browndailyherald.com/2013/10/02/creaturecast-videos-featured-ny-times-website/})
\item [\rm 2013] Article about CreatureCast. \emph{DelViscio, Jeffrey. ``Creatures, Strange and Complex, in Colorful Detail''. New York Times, September 16, 2013.} (\url{http://www.nytimes.com/2013/09/17/science/creatures-strange-and-complex-in-colorful-detail.html})
\item [\rm 2013] Interviewed, and published photographs, for \emph{Maxmen, Amy. ``Transcriptomics for the Animal Kingdom''. The Scientist, July 1, 2013.}  (\url{http://www.the-scientist.com/?articles.view/articleNo/36132/title/Transcriptomics-for-the-Animal-Kingdom/})
\item [\rm 2013] Published photograph \emph{``Evolutionary enigmas''. Science News, May 18, 2013.} (\url{http://www.sciencenews.org/view/feature/id/350120/description/Evolutionary_enigmas})
\item [\rm 2013] Interviewed about the ctenophore genome. \emph{Pennisi, Elizabeth. ``Nervous System May Have Evolved Twice''. Science, January 25, 2013} (\url{http://www.sciencemag.org/content/339/6118/391.1.full})
\item [\rm 2012] A review of my book, Practical Computing for Biologists. \emph{Aiello-Lammens, M. ``Practical Computing for Biologists''. Q. Rev. Biol. 87, 372. 2012.} ((\url{http://www.jstor.org/stable/10.1086/668128}))
\item [\rm 2012] Interviewed for an article about 3D printing. \emph{Joselow, Maxine. ``MakerBot printer transforms 2-D images into plastic models''. Brown Daily Herald, October 16, 2012.} (\url{http://goo.gl/RZgRo})
\item [\rm 2012] Seven episodes from creaturecast.org were selected for the 2012 Imagine Science Film Festival (\url{http://www.imaginesciencefilms.org/festival/2012-films/})
\item [\rm 2012] Interviewed about the World Economic Forum meeting in Davos, Switzerland. \emph{Paumgarten, Nick. ``Magic Mountain''. The New Yorker, March 5, 2012.}  (\url{http://www.newyorker.com/reporting/2012/03/05/120305fa_fact_paumgarten})
\item [\rm 2012] Article about research \emph{``Decoding the family Tree'' Division of Biology and Medicine 2010--2011 Annual Report} (\url{http://biomed.brown.edu/download/Biomed-Annual-Report.pdf})
\item [\rm 2011] Article about mollusk research. \emph{Brown, Mark. ``Squid to Snail: Biologists complete mollusk evolutionary tree''. Wired UK, October 27, 2012.} (\url{http://www.wired.co.uk/news/archive/2011-10/27/mollusk-evolutionary-tree})
\item [\rm 2011] Article about mollusk research. \emph{Weinstein, Michael. ``Mollusc family tree pruned''. Brown Daily Herald, October 27, 2011} (\url{http://www.browndailyherald.com/mollusc-family-tree-pruned-1.2661151})
\item [\rm 2011] Interviewed about art and science. \emph{Kim, Carol. ``Art for the sake of science''. Brown Daily Herald, October 13, 2011.} (\url{http://www.browndailyherald.com/art-for-the-sake-of-science-1.2652480})
\item [\rm 2011] Article about craturecast.org \emph{Long, Katherine. ``Acclaimed video podcast narrates nature's quirks'' Brown Daily Herald, October 5, 2011.} (\url{http://goo.gl/0KGqJ})
\item [\rm 2011] Interviewed about timing of animal evolution. \emph{Millius, Susan. `Biology's big bang had a long fuse`''. Science News, December 31, 2011} \url{http://goo.gl/4npj9}
\item [\rm 2011] Interviewed for article about social media and careers. \emph{Gewen, Virginia. ``Social media: Self-reflection, online''. Nature, March 30, 2011.} (\url{http://www.nature.com/naturejobs/science/articles/10.1038/nj7340-667a})
\item [\rm 2011] A review of my book, Practical Computing for Biologists. \emph{Troyanskaya, Olga. ``Don't Fear the Command Line!''. Cell, March 18, 2011.} (\url{http://dx.doi.org/10.1016/j.cell.2011.02.042})
\item [\rm 2011] Video interview at the National Science Foundation for the Waterman Award. (\url{http://goo.gl/VEi9b})
\item [\rm 2011] NSF press release for Waterman Award ceremony. (\url{http://www.nsf.gov/news/news_summ.jsp?cntn_id=119530})
\item [\rm 2011] NSF press release for Waterman Award. (\url{http://www.nsf.gov/news/news_summ.jsp?cntn_id=118915})
\item [\rm 2011] Interviewed for article about slime molds. \emph{Akst, Jef. ``Unexpected slime mold complexity''. The Scientist, March 10, 2011.} (\url{http://www.the-scientist.com/?articles.view/articleNo/29586})
\item [\rm 2011] Interviewed for article about multicellularity. \emph{Akst, Jef. ``From Simple To Complex''. The Scientist, January 1, 2011.} (\url{http://www.the-scientist.com/?articles.view/articleNo/29430/})
\item [\rm 2011] News focus highlighting genomic work. \emph{Pennisi, Elizabeth. ``Tracing the Tree of Life''. Science, February 25, 2011} (\url{http://www.sciencemag.org/content/331/6020/1005.2.full})
\item [\rm 2010] Video of presentation at PopTech 2010 (\url{https://vimeo.com/16397879})
\item [\rm 2010] PopTech Quick Take Video, produced as part of the PopTech Science Fellows program (\url{http://poptech.org/popcasts/quick\_takes\_casey\_dunn})
\item [\rm 2010] Video of presentation at the Kavli Frontiers of Science Symposium (\url{https://vimeo.com/31010067})
\item [\rm 2010] Published photograph \emph{``Big Shot: Nature's Pincushion''. Brown Medicine Winter 2010} (\url{http://goo.gl/vQq1z})
\item [\rm 2010] Article about craturecast.org \emph{Cambra, Kris. ``Creature Feature''. Brown Medicine, Fall 2010} (\url{http://goo.gl/C0gXm})
\item [\rm 2010] Article about craturecast.org \emph{Urban, Lauren. ``Creature cast''. The Scientist, March 26, 2010} (\url{http://www.the-scientist.com/?articles.view/articleNo/28881/title/Creature-cast/})
\item [\rm 2010] Research featured as one of the top 100 science stories of the year. \emph{Ruvinsky, Jessica. ``Top 100 Stories of 2009 \#43: Five Big Additions to Darwin's Theory of Evolution''. Discover Magazine, January--February 2010.} (\url{http://discovermagazine.com/2010/jan-feb/44})
\item [\rm 2010] creaturecast.org feature. \emph{Visible Insight: EPSCOR Collaboration in Rhode Island 2009--2010.} \url{http://expspace.risd.edu/?research=visible-insight}
\item [\rm 2009] Published photograph. \emph{Roach, John. ``NewGreen Bombe Sea Worms Fire Glowing Blobs''. National Geographic News, August 20, 2009.} (\url{http://news.nationalgeographic.com/news/2009/08/photogalleries/worms-glowing-bombs-green-pictures/})
\item [\rm 2009] creaturecast.org featured on National Public Radio's Science Friday (\url{http://www.npr.org/templates/story/story.php?storyId=112752248})
\item [\rm 2009] Published photograph. \emph{Fountain, Henry. ``New Find in the Pacific: Worms With Glow Sticks''. New York Times, August 25, 2009.} (\url{http://goo.gl/wcfxO})
\item [\rm 2009] Research featured as one of the top 100 science stories of the year. \emph{Rice, Jocelyn. ``Top 100 Stories of 2008 \#98: You're More Like a Sponge Than a Comb Jelly''. Discover Magazine, January 2009.} \url{http://discovermagazine.com/2009/jan/098}
\item [\rm 2008] Article about animal phylogeny research. \emph{Zimmer, Carl. ``The tree of life continues to evolve''. Boston Globe, April 28, 2008.} \url{http://www.boston.com/news/science/articles/2008/04/28/tree_of_life_continues_to_evolve/}
\item [\rm 2008] Article about animal phylogeny research. \emph{Maxmen, Amy. ``Comb Jellies Take Root in a New Tree of Animal Life''. Science, News April 5, 2008.} \url{http://www.jstor.org/stable/20465384}
\item [\rm 2008] Interviewed for article about jellyfish. \emph{Dance, Amber. ``How the jellyfish got its sting''. Nature News, September 28, 2008.} (\url{http://www.nature.com/news/2008/080928/full/news.2008.1134.html})
\item [\rm 2008] Interviewed for article about Trichoplax genome. \emph{Pennisi, Elizabeth. `` 'Simple' Animal's Genome Proves Unexpectedly Complex''. Science, August 22, 2008.} (\url{http://www.sciencemag.org/content/321/5892/1028.2.full})
\item [\rm 2008] Interviewed for article agout genome sequencing. \emph{Pennisi, Elizabeth. ``Building the Tree of Life, Genome by Genome''. Science, June 27, 2008.} (\url{http://www.sciencemag.org/content/320/5884/1716.full})
\item [\rm 2005] Published photograph (in print edition of article). \emph{Mallon, Thomas. ``Do jellyfish rule the world?'' Discover Magazine, September 2007.} \url{http://discovermagazine.com/2007/sep/do-jellyfish-rule-the-world}
\item [\rm 2005] Published photograph. \emph{Braun, David. ``Blind Sea Creature Hunts With Light''. National Geographic News, July 7, 2005.} \url{http://news.nationalgeographic.com/news/2005/07/0707_050708_redlight.html}
\item [\rm 2005] Article about siphonophore study. \emph{Broad, William J. ``Cousin of the Jellyfish Spurs Fresh Theories on Seeing Red''. The New York Times, July 12, 2005.} (\url{http://www.nytimes.com/2005/07/12/science/12red.html})
\item [\rm 2005] Article about siphonophore study. \emph{Simonite, Tom. ``Jellyfish capture prey with crimson bait''. Nature News, July 7, 2005.} (\url{http://www.nature.com/news/2005/050707/full/news050704-9.html})
\item [\rm 2005] Review of siphonophore web site. \emph{Leslie, Mitch. ``The Life Gelatinous''. Science, August 12, 2005.} (\url{http://www.sciencemag.org/content/309/5737/995.5.full})

\end{entry}

} % end \iftoggle{service}

\section{Academic Honors}

\begin{entry}{60pt}

\item [\rm 2016] Blavatnik Award for Young Scientists Finalist (\url{http://blavatnikawards.org/honorees/national-finalists/})
\item [\rm 2012] Teaching with Technology Award (Brown University)
\item [\rm 2011] Alan T. Waterman Award (National Science Foundation)
\item [\rm 2010] Named the Manning Assistant Professor of Ecology and Evolutionary Biology
\item [\rm 2010] PopTech Science and Public Leadership Fellow
\item [\rm 2010] Kavli Frontiers of Science Fellow (Sponsored by U.S. National Academy of Sciences)
\iftoggle{long}{
\item [\rm 2006] John Spangler Nicholas Prize (Yale University)
\item [\rm 2000--2003] NSF Graduate Research Fellowship (National Science Foundation)
\item [\rm 2000--2002] Sterling Fellowship (Yale University)
}% end \iftoggle{long}

\end{entry}

\iftoggle{long}{
\section{Teaching}

\subsection*{Courses}

\begin{entry}{60pt}
\item [\rm 2024 Fall] Phylogenetic Biology (EEB 354, 15 students)
\item [\rm 2024 Spring] Advanced Topics in Ecology and Evolution (EEB501, 9 students)
\item [\rm 2024 Spring] Responsible Conduct of Research (EEB545, 9 students)
\item [\rm 2023 Fall] Invertebrate Zoology Lecture (EEB255, 86 students)
\item [\rm 2023 Fall] Invertebrate Zoology Lab (EEB256L, 13 students)
\item [\rm 2022 Fall] Phylogenetic Biology (EEB 354, 19 students)
\item [\rm 2022 Spring] Advanced Topics in Ecology and Evolution (EEB501, 10 students)
\item [\rm 2022 Spring] Responsible Conduct of Research (EEB545, 10 students)
\item [\rm 2021 Fall] Advanced Topics in Ecology and Evolution (EEB500, 10 students)
\item [\rm 2021 Fall] Invertebrate Zoology Lecture (EEB255, 16 students)
\item [\rm 2021 Fall] Invertebrate Zoology Lab (EEB256L, 6 students)
\item [\rm 2020 Fall] Phylogenetic Biology (EEB 354, 22 students)
\item [\rm 2020 Fall] Advanced Philosophy of Biology (EEB 821, 8 students)
\item [\rm 2019 Fall] Philosophy of Biology (EEB321, 13 students)
\item [\rm 2019 Fall] Invertebrate Zoology Lecture (EEB255, 81 students)
\item [\rm 2019 Fall] Invertebrate Zoology Lab (EEB256L, 9 students)
\item [\rm 2019 Spring] Comparative Genomics (EEB723, 11 students)
\item [\rm 2018 Fall] Invertebrate Zoology Lecture (EEB255, 4 students)
\item [\rm 2018 Fall] Invertebrate Zoology Lab  (EEB256L, 3 students)
\item [\rm 2018 Summer] Workshop on Molecular Evolution at the Marine Biological Laboratory (Woods Hole, MA)
\item [\rm 2017 Summer] Workshop on Molecular Evolution at the Marine Biological Laboratory (Woods Hole, MA)
\item [\rm 2016 Fall] Invertebrate Zoology (Biol 0410, 43 students)
\item [\rm 2016 Spring] Phylogenetic Biology (Biol 1425, 13 students)
\item [\rm 2016 Summer] Workshop on Molecular Evolution at the Marine Biological Laboratory (Woods Hole, MA)
\item [\rm 2015 Fall] Topics in Ecology and Evolutionary Biology: interacting with data (Biol 2430, 19 students). \url{https://bitbucket.org/caseywdunn/data_interaction}
\item [\rm 2015 Fall] Invertebrate Zoology (Biol 0410, 37 students)
\item [\rm 2015 Summer] Practical Computing for Biologists Workshop, Friday Harbor Labs, University of Washington
\item [\rm 2015 Summer] Workshop on Molecular Evolution at the Marine Biological Laboratory (Woods Hole, MA)
\item [\rm 2014 Summer] Workshop on Molecular Evolution at the Marine Biological Laboratory (Woods Hole, MA)
\item [\rm 2013 Fall] Invertebrate Zoology (Biol 0410, 40 students)
\item [\rm 2013 Summer] Workshop on Molecular Evolution at the Marine Biological Laboratory (Woods Hole, MA)
\item [\rm 2013 Spring] Phylogenetic Biology (Biol 1425, 11 students)
\item [\rm 2012 Fall] Invertebrate Zoology (Biol 0410, 37 students)
\item [\rm 2012 Summer] Workshop on Molecular Evolution at the Marine Biological Laboratory (Woods Hole, MA); Guest lecture in: Bioinformatics: Biodiscovery By Computer at Summer@Brown (a high school summer course)
\item [\rm 2012 Spring] EEB IGERT Course in Reverse Ecology (3 students); Guest lectures in: Analysis of Development (Biol 1310), Topics in Science Communications: Science Journalism Practicum (Biol 0950B)
\item [\rm 2011 Fall] Invertebrate Zoology (Biol 0410, 34 students)
\item [\rm 2011 Summer] Workshop on Molecular Evolution at the Marine Biological Laboratory (Woods Hole, MA)
\item [\rm 2011 Spring] Practical Computing for Biologists NESCENT Academy at North Carolina State University (co-taught with Steve Haddock)
\item [\rm 2010 Fall] Invertebrate Zoology (Biol 0410, 33 students)
\item [\rm 2010 Spring] Origins of Multicellularity and the Evolution of Germ Line (Biol 1940Y-S01, 14 students), Directed Research/Independent Study (Biol 195, 1 students)
\item [\rm 2009 Fall] Invertebrate Zoology (Biol 0410, 18 students), Directed Research/Independent Study (Biol 195, 2 students)
\item [\rm 2009 Spring] Topics in Ecology and Evolutionary Biology (Biol 0244, 16 students, included field trip to Belize), Directed Research/Independent Study (Biol 195, 1 student)
\item [\rm 2008 Fall] Topics in Ecology and Evolutionary Biology (Biol 0243, 16 students), Directed Research/Independent Study (Biol 195, 2 student)
\item [\rm 2008 Spring] Directed Research/Independent Study (Biol 0195, 2 students)
\item [\rm 2007 Fall] Directed Research/Independent Study (Biol 0195, 1 students)

\end{entry}


\subsection*{Advising--Postdoctoral researchers}

\begin{entry}{60pt}
\item [\rm 2021--present] Sam Church
\item [\rm 2020--present] Natasha Picciani de Souza
\item [\rm 2018--2019] Yuanning Li
\item [\rm 2016--2020] Zachary Lewis
\item [\rm 2012--2016] Felipe Zapata
\item [\rm 2010--2011] Stephen Smith
\item [\rm 2009--2015] Stefan Siebert


\end{entry}


\subsection*{Advising--Graduate students (primary advisor)}

\begin{entry}{60pt}
\item [\rm 2023--present] Dalila Destanovic
\item [\rm 2021--present] Namrata Ahuja
\item [\rm 2019--present] Lauren Mellenthin
\item [\rm 2018--present] Jasmine Mah
\item [\rm 2015--2021] Alejandro Damian Serrano
\item [\rm 2015--2016] Cat Luria (co-advised with Jeremy Rich)
\item [\rm 2013--2018] August Guang (Applied Math, co-advised with Chip Lawrence)
\item [\rm 2013--2018] Catriona Munro
\item [\rm 2008--2015] Rebecca Helm


\end{entry}


\subsection*{Advising--Graduate student thesis committee member}

\begin{entry}{60pt}
\item [\rm 2025--present] Lan Wei (Yale, EEB)
\item [\rm 2025--present] Dominica Cao (Yale, MCDB)
\item [\rm 2024--present] Jacob Mathai (Yale, Immunobiology)
\item [\rm 2023--present] Roy Zang (Yale, MCDB)
\item [\rm 2023--present] Oluwatobi Oso (Yale, EEB)
\item [\rm 2023--present] Arianna Lord (Harvard, OEB)
\item [\rm 2023--present] Haysun Choi (Yale, EEB)
\item [\rm 2021--present] Delaney Farris (Yale, Genetics)
\item [\rm 2019--present] Sara Siwiecki
\item [\rm 2019--present] Pranav Kantroo (Yale, Computational Biology and Bioinformatics)
\item [\rm 2019--2024] Liam Taylor (Yale, EEB)
\item [\rm 2019--2024] Daniel Stadtmauer (Yale, EEB)
\item [\rm 2019--2021] Cecelia Harold (Yale, Genetics)
\item [\rm 2020--2023] Cory Berger(MIT/ WHOI)
\item [\rm 2019--2023] Chantal Parker (Yale, EEB)
\item [\rm 2019--2023] Ian Gilman (Yale, EEB)
\item [\rm 2018--2021] Jack Shaw (Yale, Geology)
\item [\rm 2019--2021] Nicolas Mongiardino Koch (Yale, Geology)
\item [\rm 2017--2020] Evelyn Pless (Yale, EEB)
\item [\rm 2017--2021] Sam Church (Harvard, OEB)
\item [\rm 2015--2021] Morgan Moeglein (Yale, EEB)
\item [\rm 2016--2019] Joaquin Nunez (Brown University, EEB)
\item [\rm 2019--2019] Ho Kwan Tang (Yale, EEB)
\item [\rm 2016--2017] Aislinn Rowan (Brown University, MPP)
\item [\rm 2015--2019] Christy Rhine (Brown University, MCB)
\item [\rm 2015--2019] Robert Lamb (Brown University, EEB)
\item [\rm 2015--2015] Kara Pellowe-Wagstaff (Brown University, EEB)
\item [\rm 2013--2014] Warren Francis (University of California at Santa Cruz)
\item [\rm 2012--2017] Tara Fresques (Brown University, MCB)
\item [\rm 2011--2012] Dereck Stefanik (Boston University)
\item [\rm 2011--2012] Anna Ritz (Brown University, CS)
\item [\rm 2011--2015] Ariel Camp (Brown University, EEB)
\item [\rm 2010--2012] Arturo Ruiz Villanueva (Institute of Ecology, Xalapa, Mexico)
\item [\rm 2010--2015] Cat Luria (Brown University, EEB)
\item [\rm 2009--2014] Adrian Reich (Brown University, MCB)
\item [\rm 2009--2014] Ben Ewen-Campen (Harvard University, OEB)
\item [\rm 2009--2011] John Cumbers (Brown University, MCB)
\item [\rm 2009--2015] Steven Swartz (Brown University, MCB)
\item [\rm 2009--2013] Henry Astley (Brown University, EEB)
\item [\rm 2008--2012] Matt Ogburn (Brown University, EEB)
\item [\rm 2007--2009] James Palardy (Brown University, EEB)

\end{entry}


\subsection*{Advising--Rotating graduate students}

\begin{entry}{60pt}
\item [\rm 2019] Ava Ghezelayagh
\item [\rm 2010] Matthew Booker
\item [\rm 2009] Steven Swartz


\end{entry}


\subsection*{Advising--Undergraduate students}

\begin{entry}{60pt}
\item [\rm 2019--2020] Tarek Ziad
\item [\rm 2013--2013] Rachel  Kaplan (UTRA)
\item [\rm 2013--2013] Hannah Kerman (EPSCoR Surf)
\item [\rm 2012--2015] Sam Church (EPSCoR)
\item [\rm 2012--2014] Pathikrit Bhattacharyya (EPSCoR, Royce Fellow)
\item [\rm 2012--2014] Jessica Eason (EPSCoR)
\item [\rm 2012--2013] Served as a freshman advisor for joint Brown/ RISD students
\item [\rm 2012] Robert Sandler (Royce Fellow)
\item [\rm 2012] James Weis (Thesis reader)
\item [\rm 2011--2012] Natividad Chen (Science Center Fellow, UTRA)
\item [\rm 2011--2012] Norian Caporale-Berkowitz (Beckman Fellow)
\item [\rm 2011] Cecelia Bahamon (Thesis reader)
\item [\rm 2010] Stephanie Spielman (Thesis reader)
\item [\rm 2007--2010] Orla O'Brien (Thesis advisor)
\item [\rm 2007--2010] Caitlin Feehery (participant on several projects)
\item [\rm 2007--2010] Sophia Tintori (Thesis advisor)

\end{entry}

} % end \iftoggle{long}




\bigskip

% Footer
\begin{center}
  \begin{footnotesize}
    Last updated: \today \\
  \end{footnotesize}
\end{center}

\end{document}
